\subsubsection{Range Accrual Leg Data}
\label{ss:rangeaccruallegdata}

A Range Accrual leg pays a coupon that accrues proportionally to the number of days an underlying Ibor index fixes
within a specified range $[L, U]$. For each accrual period the coupon amount is:

$$\text{Amount} = N \cdot r \cdot \frac{n}{N_{\text{obs}}} \cdot \tau$$

where $N$ is the notional, $r$ is either a fixed coupon rate or a floating rate (gearing $\times$ Ibor $+$ spread),
$n$ is the number of observation days where $L \leq \text{fixing} \leq U$, $N_{\text{obs}}$ is the total number of
observation days in the period, and $\tau$ is the day count fraction.

The \lstinline!RangeAccrualLegData! trade component node is used within the \lstinline!LegData! trade component when the
\lstinline!LegType! element is set to \emph{RangeAccrual}. It contains a nested \lstinline!FloatingLegData! node that
defines the underlying Ibor index, fixing days, gearings and spreads, plus additional elements for the range boundaries
and an optional fixed coupon rate.

Two modes are supported:
\begin{itemize}
\item \textbf{Fixed-rate mode}: When \lstinline!Coupons! are provided (non-zero), the payout per period is
  $N \cdot c \cdot (n / N_{\text{obs}}) \cdot \tau$ where $c$ is the fixed coupon rate. In this mode, gearings and
  spreads from the nested \lstinline!FloatingLegData! are ignored.
\item \textbf{Floating mode}: When \lstinline!Coupons! are omitted or set to zero, the payout per period is
  $N \cdot (g \cdot f + s) \cdot (n / N_{\text{obs}}) \cdot \tau$ where $g$ is the gearing, $f$ the index fixing and
  $s$ the spread from the nested \lstinline!FloatingLegData!.
\end{itemize}

Listing \ref{lst:rangeaccruallegdata} shows an example for a leg of type Range Accrual in fixed-rate mode.

\begin{listing}[H]
%\hrule\medskip
\begin{minted}[fontsize=\footnotesize]{xml}
      <LegData>
        <LegType>RangeAccrual</LegType>
        <Payer>false</Payer>
        <Currency>EUR</Currency>
        <Notionals>
          <Notional>10000000</Notional>
        </Notionals>
        <DayCounter>ACT/360</DayCounter>
        <PaymentConvention>ModifiedFollowing</PaymentConvention>
        <ScheduleData>
          <Rules>
            <StartDate>2024-01-15</StartDate>
            <EndDate>2029-01-15</EndDate>
            <Tenor>6M</Tenor>
            <Calendar>TARGET</Calendar>
            <Convention>ModifiedFollowing</Convention>
            <TermConvention>ModifiedFollowing</TermConvention>
            <Rule>Forward</Rule>
          </Rules>
        </ScheduleData>
        <RangeAccrualLegData>
          <FloatingLegData>
            <Index>EUR-EURIBOR-6M</Index>
            <FixingDays>2</FixingDays>
          </FloatingLegData>
          <Coupons>
            <Coupon>0.05</Coupon>
          </Coupons>
          <LowerBounds>
            <LowerBound>0.02</LowerBound>
          </LowerBounds>
          <UpperBounds>
            <UpperBound>0.06</UpperBound>
          </UpperBounds>
        </RangeAccrualLegData>
      </LegData>
\end{minted}
\caption{Range Accrual leg data (fixed-rate mode)}
\label{lst:rangeaccruallegdata}
\end{listing}

Listing \ref{lst:rangeaccruallegdata_floating} shows an example in floating mode with gearing and spread.

\begin{listing}[H]
%\hrule\medskip
\begin{minted}[fontsize=\footnotesize]{xml}
        <RangeAccrualLegData>
          <FloatingLegData>
            <Index>EUR-EURIBOR-6M</Index>
            <FixingDays>2</FixingDays>
            <Spreads>
              <Spread>0.005</Spread>
            </Spreads>
            <Gearings>
              <Gearing>1.5</Gearing>
            </Gearings>
          </FloatingLegData>
          <LowerBounds>
            <LowerBound>0.01</LowerBound>
          </LowerBounds>
          <UpperBounds>
            <UpperBound>0.05</UpperBound>
          </UpperBounds>
        </RangeAccrualLegData>
\end{minted}
\caption{Range Accrual leg data (floating mode)}
\label{lst:rangeaccruallegdata_floating}
\end{listing}

The \lstinline!RangeAccrualLegData! block contains the following elements:

\begin{itemize}

\item FloatingLegData: A nested \lstinline!FloatingLegData! node (see \ref{ss:floatingleg_data}) that defines the
  underlying Ibor index used for range observations and, in floating mode, for coupon rate determination. Only Ibor
  indices are supported. The relevant sub-elements are:

  \begin{itemize}
  \item Index: The Ibor index whose daily fixings are observed against the range bounds, e.g.\ \emph{EUR-EURIBOR-6M}.
    See Table \ref{tab:indices} for allowable values.
  \item FixingDays [Optional]: The number of business days before the observation date to take the index fixing.
    Defaults to the index's fixing days if omitted.
  \item Spreads [Optional]: Spread(s) added to the index fixing in floating mode. Ignored in fixed-rate mode. Expressed
    in decimal form, e.g.\ \emph{0.005} is 50 bp.
  \item Gearings [Optional]: Gearing factor(s) applied to the index fixing in floating mode. Ignored in fixed-rate
    mode. Defaults to \emph{1} if omitted.
  \end{itemize}

\item Coupons [Optional]: This node contains child elements of type \lstinline!Coupon! specifying the fixed coupon
  rate(s) for fixed-rate mode. If this is a single value, it applies to all periods. A vector of values or a dated
  vector of values (using \lstinline!startDate! attributes) can be used for varying rates.

  When \lstinline!Coupons! are provided with non-zero values, the leg operates in fixed-rate mode and the gearings and
  spreads from \lstinline!FloatingLegData! are ignored.

  If the entire \lstinline!Coupons! section is omitted or all values are zero, the leg operates in floating mode.

  Allowable values: Each child element can take any real number. The rate is expressed in decimal form, e.g.\
  \emph{0.05} is a rate of 5\%.

\item LowerBounds: This node contains child elements of type \lstinline!LowerBound! specifying the lower bound(s) of
  the accrual range. The observation index fixing must be greater than or equal to the lower bound for a day to count
  towards the accrual. This can be a single value, a vector of values, or a dated vector of values (using
  \lstinline!startDate! attributes).

  Allowable values: Each child element can take any real number. The bound is expressed in decimal form, e.g.\
  \emph{0.02} is a bound of 2\%.

\item UpperBounds: This node contains child elements of type \lstinline!UpperBound! specifying the upper bound(s) of
  the accrual range. The observation index fixing must be less than or equal to the upper bound for a day to count
  towards the accrual. This can be a single value, a vector of values, or a dated vector of values (using
  \lstinline!startDate! attributes).

  Allowable values: Each child element can take any real number. The bound is expressed in decimal form, e.g.\
  \emph{0.06} is a bound of 6\%.

\end{itemize}
